%% Question bank chapter file - generated by bank_tools.py build (do not edit while the bank is still being built:
%% it is regenerated from the per-test data in Question Bank\data). Every question is one bankq block:
%%   \begin{bankq}[id=..., source={...}, paper=.., qno=.., type=..., subtopic={...}, answer={...}, answer_src=key|solved,
%%                 working={...}, difficulty=1-10, diff_parts={K S I C F}, diff_note={...},
%%                 group=<passage id or empty>, also_in={duplicate copies in other tests}]
%% difficulty = min(10, K+S+I+C+F-1): K concepts 1-3, S steps 1-4, I insight 0-2, C computation 0-2, F format 0-1
%%   ...question LaTeX, options, figures (\bfig / \bfigi -> img/)...
%%   \end{bankq}
%% Passages shared by several questions are bankpassage blocks; their questions carry group=<passage id>.
\documentclass[10pt]{article}
\usepackage[a4paper,margin=1.8cm]{geometry}
\usepackage[T1]{fontenc}
\usepackage{amsmath,amssymb,mathtools}
\usepackage{graphicx}
\usepackage[export]{adjustbox}
\usepackage{xcolor,needspace,keyval}
\usepackage{enumitem}
\usepackage{booktabs,array,tabularx}
\usepackage[hidelinks]{hyperref}
\setlength{\parindent}{0pt}
\setlength{\parskip}{3pt}
\makeatletter
\define@key{bq}{id}{\def\bq@id{#1}}\define@key{bq}{source}{\def\bq@src{#1}}\define@key{bq}{paper}{\def\bq@paper{#1}}
\define@key{bq}{qno}{\def\bq@qno{#1}}\define@key{bq}{type}{\def\bq@type{#1}}\define@key{bq}{subtopic}{\def\bq@sub{#1}}
\define@key{bq}{answer}{\def\bq@ans{#1}}\define@key{bq}{answer_src}{\def\bq@asrc{#1}}\define@key{bq}{working}{\def\bq@work{#1}}
\define@key{bq}{group}{\def\bq@group{#1}}\define@key{bq}{also_in}{\def\bq@also{#1}}
\define@key{bq}{difficulty}{\def\bq@diff{#1}}\define@key{bq}{diff_parts}{\def\bq@dparts{#1}}\define@key{bq}{diff_note}{\def\bq@dnote{#1}}
\newenvironment{bankq}[1][]{\def\bq@work{}\def\bq@also{}\def\bq@group{}\def\bq@diff{?}\def\bq@dparts{}\def\bq@dnote{}\setkeys{bq}{#1}%
  \par\needspace{4\baselineskip}\medskip\noindent{\small\textbf{\bq@id}\hfill\color{gray}[\bq@src, Paper \bq@paper, Q\bq@qno]
  \bq@type\ $\cdot$ \bq@sub\ $\cdot$ \textbf{D\bq@diff}}\par\nopagebreak\relax}
  {\par{\small\color{teal}Answer: \textbf{\bq@ans}\ifx\bq@work\@empty\else\ (solved: \bq@work)\fi
  \ifx\bq@also\@empty\else\par\color{gray}Also in: \bq@also\fi}\par}
\define@key{bp}{id}{\def\bp@id{#1}}\define@key{bp}{source}{\def\bp@src{#1}}\define@key{bp}{questions}{\def\bp@qs{#1}}
\newenvironment{bankpassage}[1][]{\setkeys{bp}{#1}\par\needspace{4\baselineskip}\bigskip\noindent
  {\small\textbf{Passage \bp@id}\hfill\color{gray}[\bp@src; questions \bp@qs]}\par\nopagebreak\relax\itshape}{\par}
\makeatother
\newcommand{\bfig}[2]{\par\smallskip{\centering\includegraphics[width=#1pt,max width=0.62\linewidth,max height=0.32\textheight]{img/#2}\par}\smallskip}
\newcommand{\bfigi}[2]{\includegraphics[width=#1pt,max width=\linewidth,max height=0.32\textheight,valign=c]{img/#2}}
\newcommand{\bmissing}[1]{\textit{[figure missing in the source conversion: \detokenize{#1}]}}
\newenvironment{qoptions}{\begin{itemize}[leftmargin=1.8em,itemsep=1pt,topsep=2pt,label={}]}{\end{itemize}}
\newcommand{\opt}[2]{\item[(#1)] #2}
\providecommand{\fwstart}{}\providecommand{\fwend}{}\providecommand{\qfullsplit}{}
\newcolumntype{L}[1]{>{\raggedright\arraybackslash}p{#1}}
\newcolumntype{C}[1]{>{\centering\arraybackslash}p{#1}}
\newcolumntype{R}[1]{>{\raggedleft\arraybackslash}p{#1}}
\providecommand{\sectionhdr}[1]{}\providecommand{\subsectionhdr}[1]{}


\begin{document}
{\Large\bfseries Properties of Matter --- Question Bank}\par\smallskip
{\small Physics chapter P23. Questions from all converted test papers, grouped by source series. Answers: printed key where it matches the paper, otherwise solved (answer\_src=solved, with working). Exact repeats across papers appear once, with also\_in listing the other copies.}\par\medskip\hrule

\section*{TOAS papers}
\begin{bankq}[id=TOAS11-P1-PHY-S3-Q02, source={TOAS - Test 11 (27 Jun 2022)}, paper=1, qno=2, type=numerical, subtopic={Young's modulus}, answer={1.00}, answer_src=key, difficulty=4, diff_parts={K1 S2 I1 C0 F1}, diff_note={concepts: elongation of a wire with varying radius | steps: radius as a linear function of position; integrate the strain; evaluate | I: integral of dx over the varying area}]
A body of mass 3.14 kg is suspended from one end of a wire of length 10 m. The radius of the wire is changing uniformly from $10\times10^{-4}$ m at one end to $5\times10^{-4}$ m at the other end. Find the change in length of the wire in mm. Young's modulus of the material of the wire is $2\times10^{11}$ N/m$^2$. (mass of wire is negligible) Take : $\pi=3.14$ and $g=10$ m/s$^2$
\bfig{94.3}{TOAS11_p1_phy_s3_q02_1.png}
\end{bankq}

\section*{TAAS papers}
\begin{bankq}[id=TAAS02-P2-PHY-S4-Q08, source={TAAS - Test 02 (30 Jul 2023)}, paper=2, qno=8, type=numerical, subtopic={Young's modulus}, answer={4}, answer_src=key, difficulty=6, diff_parts={K1 S2 I2 C1 F1}, diff_note={concepts: extension of a rod of varying cross section | steps: strain of a thin slice; integrate along the length; evaluate}]
A rod of length $\ell=2$ and varying cross section as shown is acted upon by two equal and opposite forces $F=4$N. Elastic constant of material of rod is $E=2\times10^{11}$ N/m$^2$. The increase in length of the rod is $x_0\times10^{-11}\ \ell$n(2)m. Find $x_0$. [Given, $b=1$m, $h_1=1$m and $h_2=2$m]
\bfig{161.1}{TAAS02_p2_phy_s4_q08_1.png}
\end{bankq}

\section*{TOPS papers}
\begin{bankq}[id=TOPS02-P2-PHY-S2-Q02, source={TOPS - Test 02 (28 Aug 2022)}, paper=2, qno=2, type=numerical, subtopic={Elasticity}, answer={16.67}, answer_src=solved, working={Tension in the cord is $mg=0.05$ N everywhere above the insect, so strain is $mg/(Kl_0)=2$ and each natural length stretches to $3\times$. Insect at 10 cm from the ceiling means $10/3$ cm of natural length above it; the cord below carries no tension, leaving $20/3$ cm. $L=10+20/3\approx16.67$ cm}, difficulty=8, diff_parts={K2 S3 I2 C1 F1}, diff_note={concepts: Hooke's law of a cord, strain | steps: tension equals the weight; strain per unit natural length; insect position in natural length; slack part; total length | cord stretched by its own strain}]
An elastic cord of force constant $K=0.25$ N/m$^2$ and natural length 10 cm is hanging from a ceiling. An insect of mass 5 g is sitting at the lower end of the cord and the system is in equilibrium. The insect start climbing up the spring toward the ceiling. Assume that there is no slipping and insect is moving slowly. Length of the cord, when insect is $1/3^{\text{rd}}$ of its original distance from ceiling, is $L$ cm, then $L$ is equal to (given that $g=10$ m/s$^2$)
\bfig{100.2}{TOPS02_p2_phy_s2_q02_1.png}
\end{bankq}
\begin{bankq}[id=TOPS04-P1-PHY-S4-Q07, source={TOPS - Test 04 (18 Dec 2022)}, paper=1, qno=7, type=numerical, subtopic={Young's modulus}, answer={2}, answer_src=key, difficulty=8, diff_parts={K2 S3 I2 C1 F1}, diff_note={concepts: Young's modulus, torque balance | steps: equal extension of the rods; areas and lengths; force ratio; torque balance; position}]
A very stiff bar $(AB)$ of negligible mass is suspended horizontally by two vertical rods as shown in figure. Length of the bar is $2.5L$. The steel rod has length $L$ and cross sectional radius of $r$ and the brass rod has length $2L$ and cross sectional radius of $2r$. A vertically downward force $F$ is applied to the bar at a distance $x$ from the steel rod and the bar remains horizontal. Find the value of $x$ if it is given that ratio of Young's modulus of steel and brass is $\dfrac{Y_s}{Y_B}=2$. (Given : $L=1.6$ m)
\bfig{160.9}{TOPS04_p1_phy_s4_q07_1.png}
\end{bankq}
\begin{bankq}[id=TOPS04-P1-PHY-S4-Q09, source={TOPS - Test 04 (18 Dec 2022)}, paper=1, qno=9, type=numerical, subtopic={Stress and strain}, answer={3}, answer_src=key, difficulty=7, diff_parts={K2 S2 I2 C1 F1}, diff_note={concepts: stress and strain, rotation | steps: tension at a radius; integrate the strain; total extension; convert to cm}]
A thin rod of length 3m is being rotated in horizontal plane with 10 rad/sec about its one end. The change in length (in cm) of the rod due to rotation is (Given : Young's modulus of elasticity of the wire, $Y=3\times10^8$ N/m$^2$. Density of the rod is $10^4$ kg/m$^3$)
\bfig{92.3}{TOPS04_p1_phy_s4_q09_1.png}
\end{bankq}
\begin{bankq}[id=TOPS04-P2-PHY-S2-Q06, source={TOPS - Test 04 (18 Dec 2022)}, paper=2, qno=6, type=numerical, subtopic={Young's modulus}, answer={14.00}, answer_src=key, difficulty=10, diff_parts={K3 S3 I2 C2 F1}, diff_note={concepts: Young's modulus, circular motion, extension | steps: energy to the lowest point; tension; extension; geometry to the floor; speed squared}]
A sphere of radius 10 cm and mass 25 kg is attached to the lower end of a steel wire of radius 0.05 cm and length 500 cm which is suspended from the ceiling of a room. The point of support is 521 cm above the floor. When the sphere is set swinging as a simple pendulum, its lowest point just grazes the floor, if the velocity of the ball at its lowest position (in m/s) is $v$ then find $v^2$. Given : $Y$ for steel $=20\times10^{10}$ Nm$^{-2}$, $g=9.8$ ms$^{-2}$. Round off the answer to closest integer.
\end{bankq}
\begin{bankq}[id=TOPS05-P2-PHY-S1-Q03, source={TOPS - Test 05 (14 Jan 2023)}, paper=2, qno=3, type=single, subtopic={Stress and strain}, answer={A}, answer_src=key, difficulty=6, diff_parts={K2 S2 I2 C1 F0}, diff_note={concepts: torsion, modulus of rigidity | steps: shear strain; shear stress; torque of the thin shell; integrate}]
A pipe of radius $R$ and thickness $t$ ($\ll R$) is fixed at one end as shown in the figure. Length of the pipe is $\ell$ and modulus of rigidity $\eta$. Find torque required to produce a torsion of $\theta$.
\bfig{108.3}{TOPS05_p2_phy_s1_q03_1.png}
\begin{qoptions}
\opt{A}{$\dfrac{2\pi R^3t\eta}{\ell}\theta$}
\opt{B}{$\dfrac{\pi R^3t\eta}{\ell}\theta$}
\opt{C}{$\dfrac{2\pi R^3t\eta}{3\ell}\theta$}
\opt{D}{$\dfrac{\pi R^3t\eta}{4\ell}\theta$}
\end{qoptions}
\end{bankq}
\begin{bankq}[id=TOPS05-P2-PHY-S2-Q05, source={TOPS - Test 05 (14 Jan 2023)}, paper=2, qno=5, type=multi, subtopic={Stress and strain}, answer={A,B,C}, answer_src=key, difficulty=6, diff_parts={K2 S2 I1 C1 F1}, diff_note={concepts: Poisson's ratio, volume strain | steps: lateral strain; volume change relation; check each}]
Choose the correct statement for Poisson's ratio of a body as shown in the figure.
\bfig{108.3}{TOPS05_p2_phy_s2_q05_1.png}
\begin{qoptions}
\opt{A}{If $\sigma=0.5$, volume remain constant}
\opt{B}{If $\sigma<0.5$, volume increases}
\opt{C}{if $\sigma=0.2$ and axial strain is 2\%, percentage change in volume is 1.2\%.}
\opt{D}{None of these}
\end{qoptions}
\end{bankq}

\section*{TMAS papers}
\begin{bankq}[id=TMAS03-P1-PHY-S1-Q04, source={TMAS - Test 03 (10 Sep 2023)}, paper=1, qno=4, type=multi, subtopic={Stress and strain}, answer={B,C}, answer_src=solved, working={$T=(M/L)\omega^2R^2$ with $R=L/2\pi$ gives $T=\dfrac{M\omega^2L}{4\pi^2}$; stress $T/A$; elongation $=\dfrac{T}{AY}L$}, difficulty=6, diff_parts={K2 S2 I1 C1 F1}, diff_note={concepts: rotating ring, elasticity | steps: tension from the centripetal force; stress; strain; elongation}]
A wire of mass $M$, length $L$ and cross-section $A$ is shaped into a ring. It is placed on smooth horizontal surface and is rotated about its axis with constant angular speed $\omega$ as shown. If Young's modulus of material of wire is $Y$, then :-
\bfig{108.3}{TMAS03_p1_phy_s1_q04_1.png}
\begin{qoptions}
\opt{A}{Tension developed in the wire is $\dfrac{M\omega^2L}{\pi^2}$.}
\opt{B}{Elongation produced in the wire is $\dfrac{M\omega^2L^2}{4\pi^2YA}$.}
\opt{C}{Stress developed in the wire is $\dfrac{M\omega^2L}{4\pi^2A}$.}
\opt{D}{Strain developed in the wire is $\dfrac{M\omega^2L}{\pi^2YA}$.}
\end{qoptions}
\end{bankq}

\section*{TNAS papers}
\begin{bankq}[id=TNAS01-P1-PHY-S2-Q07, source={TNAS - Test 01 (05 Feb 2023)}, paper=1, qno=7, type=numerical, subtopic={Young's modulus}, answer={2}, answer_src=key, difficulty=6, diff_parts={K1 S2 I2 C1 F1}, diff_note={concepts: equilibrium and elongation of wires | steps: tensions in the wires; extension of each; ratio}]
A block of weight $200$ N is suspended by copper and steel wires of same cross sectional area $0.5$ cm$^2$ and, length $\dfrac{\sqrt3}{2}$m and $\dfrac12$m, respectively. Their other ends are fixed on a ceiling as shown in figure. The angles subtended by copper and steel wires with ceiling are $30^\circ$ and $60^\circ$, respectively. If elongation in copper wire is ($\Delta\ell_1$) and elongation in steel wire is ($\Delta\ell_2$), then the ratio $\dfrac{\Delta\ell_1}{\Delta\ell_2}$ is \underline{\hspace{0.8cm}}.

[Young's modulus for copper and steel are $1\times10^{11}$ N/m$^2$ and $2\times10^{11}$ N/m$^2$, respectively.]
\bfig{124.1}{TNAS01_p1_phy_s2_q07_1.png}
\end{bankq}
\begin{bankq}[id=TNAS04-P2-PHY-S3-Q05, source={TNAS - Test 04 (18 Jun 2023)}, paper=2, qno=5, type=numerical, subtopic={Stress and strain}, answer={19.00}, answer_src=key, difficulty=6, diff_parts={K1 S2 I2 C1 F1}, diff_note={concepts: elastic energy in a cube under its own weight | steps: stress at depth; energy density; integrate over the thirds; ratio}]
A uniform cube is kept on a horizontal surface. What is the ratio of strain energy stored in lower one third to the upper one third of cube.
\end{bankq}
\end{document}
